\subsection{Double Duals and Dual Maps}
Throughout this subsection, all vector spaces are over the same fixed field $\mathbb{F}$.

\begin{definition} \label{def:algebraic-double-dual-and-natural-map} \hlt{Algebraic Double Dual and Natural Map}
\begin{enumerate}[label=(\roman*)]
\setlength{\itemsep}{0pt}
\item The \hlt{algebraic double dual} of $\mathcal{V}$ is $\mathcal{V}''\defeq(\mathcal{V}')'$.
\item The \hlt{natural evaluation map} $\Lambda_{\mathcal{V}}:\mathcal{V}\to\mathcal{V}''$ is defined by $(\forall\,\mathbf{v}\in\mathcal{V})(\forall\,\varphi\in\mathcal{V}')\,\Lambda_{\mathcal{V}}(\mathbf{v})(\varphi)\defeq\varphi(\mathbf{v})$.\\
For $\mathbf{v}\in\mathcal{V}$, $(\forall\,a,b\in\mathbb{F})(\forall\,\varphi,\psi\in\mathcal{V}')\,\Lambda_{\mathcal{V}}(\mathbf{v})(a\varphi+b\psi)=a\varphi(\mathbf{v})+b\psi(\mathbf{v})$ by pointwise operations in Proposition \ref{prop:linear-maps-form-vector-space}, so $\Lambda_{\mathcal{V}}(\mathbf{v})\in\mathcal{V}''$.
\item The space $\mathcal{V}$ is \hlt{algebraically reflexive} iff $\Lambda_{\mathcal{V}}$ is surjective.
\end{enumerate}
\end{definition}

\begin{theorem} \label{thm:algebraic-reflexivity-natural-map} \hlt{Algebraic Reflexivity}\\
Let $\mathcal{V}$ be a vector space. Claims (iii)--(v) assume Zorn's lemma, and (v) additionally assumes the Choice Principle \ref{princ:axiom-of-choice-analysis}. Then:
\begin{enumerate}[label=(\roman*)]
\setlength{\itemsep}{0pt}
\item $\Lambda_{\mathcal{V}}\in\LL(\mathcal{V},\mathcal{V}'')$;
\item if $\dim(\mathcal{V})<\infty$, then $\Lambda_{\mathcal{V}}$ is an isomorphism;
\item $\Lambda_{\mathcal{V}}$ is injective;
\item $\Lambda_{\mathcal{V}}$ is an isomorphism iff $\mathcal{V}$ is algebraically reflexive;
\item $\mathcal{V}$ is algebraically reflexive iff $\dim(\mathcal{V})<\infty$.
\end{enumerate}
\begin{proof}
\begin{enumerate}[label=(\roman*)]
\setlength{\itemsep}{0pt}
\item[$(i)$] For $a,b\in\mathbb{F}$, $\mathbf{u},\mathbf{v}\in\mathcal{V}$, and $\varphi\in\mathcal{V}'$, one has $\Lambda_{\mathcal{V}}(a\mathbf{u}+b\mathbf{v})(\varphi)=\varphi(a\mathbf{u}+b\mathbf{v})=\bigl(a\Lambda_{\mathcal{V}}(\mathbf{u})+b\Lambda_{\mathcal{V}}(\mathbf{v})\bigr)(\varphi)$.\\
Hence $\Lambda_{\mathcal{V}}(a\mathbf{u}+b\mathbf{v})=a\Lambda_{\mathcal{V}}(\mathbf{u})+b\Lambda_{\mathcal{V}}(\mathbf{v})$ by Assumption \ref{assump:function-extensionality}, so $\Lambda_{\mathcal{V}}\in\LL(\mathcal{V},\mathcal{V}'')$.
\item[$(ii)$] Suppose $\dim(\mathcal{V})<\infty$, let $\beta=(\mathbf{v}_1,\ldots,\mathbf{v}_n)$ be a basis, and let $\beta'=(\varphi_1,\ldots,\varphi_n)$ be its dual basis.\\
If $\mathbf{v}=\sum\limits_{j=1}^{n}c_j\mathbf{v}_j\in\ker(\Lambda_{\mathcal{V}})$, then $(\forall\,k\in\{1,\ldots,n\})\,c_k=\Lambda_{\mathcal{V}}(\mathbf{v})(\varphi_k)=0$; hence $\mathbf{v}=\mathbf{0}$, also for $n=0$. Thus $\Lambda_{\mathcal{V}}$ is injective by Proposition \ref{prop:injective-iff-kernel-zero}. By Theorem \ref{thm:dual-basis-and-dimension}(ii) for $\mathcal{V}$, $\dim(\mathcal{V}')=\dim(\mathcal{V})<\infty$; by the same theorem for $\mathcal{V}'$, $\dim(\mathcal{V}'')=\dim(\mathcal{V}')$.\\
Therefore $\Lambda_{\mathcal{V}}$ is surjective by Corollary \ref{cor:equal-dimension-kernel-range}(ii) and is an isomorphism by Proposition \ref{prop:linear-map-invertible-iff-bijective}.
\item[$(iii)$] Assume Zorn's lemma. By Proposition \ref{prop:injective-iff-kernel-zero}, it suffices to prove $\ker(\Lambda_{\mathcal{V}})=\{\mathbf{0}\}$.\\
If $\mathbf{v}\in\ker(\Lambda_{\mathcal{V}})$, then $(\forall\,\varphi\in\mathcal{V}')\,\varphi(\mathbf{v})=0$, so $\mathbf{v}=\mathbf{0}$ by Proposition \ref{prop:functionals-separate-points}(ii).
\item[$(iv)$] ($\Rightarrow$) If $\Lambda_{\mathcal{V}}$ is an isomorphism, then it is surjective by Proposition \ref{prop:linear-map-invertible-iff-bijective}; hence $\mathcal{V}$ is algebraically reflexive by Definition \ref{def:algebraic-double-dual-and-natural-map}(iii).
\item[$(iv)$] ($\Leftarrow$) By Proposition \ref{prop:linear-map-invertible-iff-bijective}, it suffices to prove that $\Lambda_{\mathcal{V}}$ is bijective. It is surjective by algebraic reflexivity and injective by (iii).
\item[$(v)$] ($\Rightarrow$) Suppose $\mathcal{V}$ is algebraically reflexive. Then $\Lambda_{\mathcal{V}}$ is an isomorphism by (iv).\\
If $\dim(\mathcal{V})$ were infinite, then $\dim(\mathcal{V})<\dim(\mathcal{V}')<\dim(\mathcal{V}'')$ by Theorem \ref{thm:dimension-of-algebraic-dual}(i),(ii), first for $\mathcal{V}$ and then for $\mathcal{V}'$.\\
Choose a Hamel basis $\mathcal{B}$ of $\mathcal{V}$ by Theorem \ref{thm:hamel-basis-extension}(ii). The set $\Lambda_{\mathcal{V}}(\mathcal{B})$ is a Hamel basis of $\mathcal{V}''$ by linearity, injectivity, and surjectivity of $\Lambda_{\mathcal{V}}$: linearity and injectivity preserve finite independence, while linearity and surjectivity preserve spanning.\\
Thus $\dim(\mathcal{V})=\dim(\mathcal{V}'')$ by Definition \ref{def:hamel-dimension}, a contradiction. Hence $\dim(\mathcal{V})<\infty$.
\item[$(v)$] ($\Leftarrow$) If $\dim(\mathcal{V})<\infty$, then $\Lambda_{\mathcal{V}}$ is surjective by (ii) and Proposition \ref{prop:linear-map-invertible-iff-bijective}; hence $\mathcal{V}$ is algebraically reflexive by Definition \ref{def:algebraic-double-dual-and-natural-map}(iii).
\end{enumerate}
\hlt{Structural Commentary}: Evaluation proves linearity, separation proves injectivity, and dual dimensions obstruct infinite reflexivity.\\
\hlt{Key Insight}: Canonical evaluation is surjective exactly before dual dimension grows.
\end{proof}
\end{theorem}

\begin{perspective} \hlt{Why the Double Dual Is Natural}\\
An isomorphism $\mathcal{V}\cong\mathcal{V}'$ for finite-dimensional $\mathcal{V}$ generally depends on a chosen basis.\
In contrast, $\Lambda_{\mathcal{V}}:\mathcal{V}\to\mathcal{V}''$ is defined solely by evaluation and requires no basis.
\end{perspective}

\begin{definition} \label{def:dual-map} \hlt{Dual Map}\\
If $T\in\LL(\mathcal{V},\mathcal{W})$, its \hlt{dual map} is $T':\mathcal{W}'\to\mathcal{V}'$ defined by $(\forall\,\varphi\in\mathcal{W}')\,T'\varphi\defeq\varphi T$.\\
The relation $(\forall\,\varphi\in\mathcal{W}')\,\varphi T\in\LL(\mathcal{V},\mathbb{F})=\mathcal{V}'$ holds by Proposition \ref{prop:linear-map-product-algebra}(i). Equivalently, $(\forall\,\varphi\in\mathcal{W}')(\forall\,\mathbf{v}\in\mathcal{V})\,(T'\varphi)(\mathbf{v})=\varphi(T\mathbf{v})$.
\end{definition}

\begin{proposition} \label{prop:dual-map-is-linear} \hlt{Dual Maps Are Linear}\\
If $T\in\LL(\mathcal{V},\mathcal{W})$, then $T'\in\LL(\mathcal{W}',\mathcal{V}')$.
\begin{proof}
For $a,b\in\mathbb{F}$, $\varphi,\psi\in\mathcal{W}'$, and $\mathbf{v}\in\mathcal{V}$, one has $\bigl(T'(a\varphi+b\psi)\bigr)(\mathbf{v})=(a\varphi+b\psi)(T\mathbf{v})=\bigl(aT'\varphi+bT'\psi\bigr)(\mathbf{v})$.\\
Hence $T'(a\varphi+b\psi)=aT'\varphi+bT'\psi$ by Assumption \ref{assump:function-extensionality}, so $T'\in\LL(\mathcal{W}',\mathcal{V}')$ by Definition \ref{def:linear-map}.\\
\hlt{Structural Commentary}: Pointwise evaluation reduces linearity of the pullback to linearity of each functional.\\
\hlt{Key Insight}: Precomposition preserves linear combinations while reversing map direction.
\end{proof}
\end{proposition}

\begin{theorem} \label{thm:algebra-of-dual-maps} \hlt{Algebra of Dual Maps}
\begin{enumerate}[label=(\roman*)]
\setlength{\itemsep}{0pt}
\item If $S,T\in\LL(\mathcal{V},\mathcal{W})$ and $a,b\in\mathbb{F}$, then $(aS+bT)'=aS'+bT'$.
\item If $T\in\LL(\mathcal{U},\mathcal{V})$ and $S\in\LL(\mathcal{V},\mathcal{W})$, then $(ST)'=T'S'$.
\item $(I_{\mathcal{V}})'=I_{\mathcal{V}'}$.
\end{enumerate}
\begin{proof}
\begin{enumerate}[label=(\roman*)]
\setlength{\itemsep}{0pt}
\item For $\varphi\in\mathcal{W}'$ and $\mathbf{v}\in\mathcal{V}$, $\bigl((aS+bT)'\varphi\bigr)(\mathbf{v})=\varphi(aS\mathbf{v}+bT\mathbf{v})=\bigl(aS'\varphi+bT'\varphi\bigr)(\mathbf{v})$.\\
Hence $(aS+bT)'=aS'+bT'$ by Assumption \ref{assump:function-extensionality} on $\mathcal{V}$ and $\mathcal{W}'$.
\item For $\varphi\in\mathcal{W}'$ and $\mathbf{u}\in\mathcal{U}$, $\bigl((ST)'\varphi\bigr)(\mathbf{u})=\varphi(S(T\mathbf{u}))=\bigl(T'(S'\varphi)\bigr)(\mathbf{u})$. Hence $(ST)'=T'S'$ by Assumption \ref{assump:function-extensionality} on $\mathcal{U}$ and $\mathcal{W}'$.
\item For $\varphi\in\mathcal{V}'$ and $\mathbf{v}\in\mathcal{V}$, $\bigl((I_{\mathcal{V}})'\varphi\bigr)(\mathbf{v})=\varphi(I_{\mathcal{V}}\mathbf{v})=\varphi(\mathbf{v})$. Hence $(I_{\mathcal{V}})'=I_{\mathcal{V}'}$ by Assumption \ref{assump:function-extensionality} on $\mathcal{V}$ and $\mathcal{V}'$.
\end{enumerate}
\hlt{Structural Commentary}: Function extensionality lifts pointwise linearity, composition, and identity calculations to dual maps.\\
\hlt{Key Insight}: Precomposition reverses products but preserves linear structure.
\end{proof}
\end{theorem}

\begin{corollary} \label{cor:dual-map-inverse} \hlt{Dual Map of an Inverse}\\
If $T\in\LL(\mathcal{V},\mathcal{W})$ is an isomorphism, then:
\begin{enumerate}[label=(\roman*)]
\setlength{\itemsep}{0pt}
\item $T'$ is an isomorphism;
\item $(T^{-1})'=(T')^{-1}$.
\end{enumerate}
\begin{proof}
Set $S\defeq(T^{-1})'$. Then $S\in\LL(\mathcal{V}',\mathcal{W}')$ by Proposition \ref{prop:dual-map-is-linear}.
\begin{enumerate}[label=(\roman*)]
\setlength{\itemsep}{0pt}
\item By Theorem \ref{thm:algebra-of-dual-maps}(ii)--(iii), $T'S=(T^{-1}T)'=(I_{\mathcal{V}})'=I_{\mathcal{V}'}$ and $ST'=(TT^{-1})'=(I_{\mathcal{W}})'=I_{\mathcal{W}'}$.\\
Thus $S$ is a two-sided inverse of $T'$, and $T'$ is an isomorphism by Definition \ref{def:invertible-linear-map-and-isomorphism}.
\item The identity $S=(T')^{-1}$ holds by Proposition \ref{prop:inverse-linear-map-is-unique}; hence $(T^{-1})'=(T')^{-1}$.
\end{enumerate}
\hlt{Structural Commentary}: Reversed product identities construct the dual inverse; uniqueness identifies it.\\
\hlt{Key Insight}: Contravariance turns both inverse identities into dual identities.
\end{proof}
\end{corollary}

\begin{corollary} \label{cor:double-dual-map-identifies-with-original-map} \hlt{Double Dual Maps Identify with Original Maps}\\
Let $T\in\LL(\mathcal{V},\mathcal{W})$ and set $T''\defeq(T')'$. Then:
\begin{enumerate}[label=(\roman*)]
\setlength{\itemsep}{0pt}
\item $T''\Lambda_{\mathcal{V}}=\Lambda_{\mathcal{W}}T$;
\item if $\dim(\mathcal{V}),\dim(\mathcal{W})<\infty$, then $T''=\Lambda_{\mathcal{W}}T\Lambda_{\mathcal{V}}^{-1}$.
\end{enumerate}
\begin{proof}
$T'\in\LL(\mathcal{W}',\mathcal{V}')$ and $T''\in\LL(\mathcal{V}'',\mathcal{W}'')$ by Proposition \ref{prop:dual-map-is-linear}, first for $T$ and then for $T'$.
\begin{enumerate}[label=(\roman*)]
\setlength{\itemsep}{0pt}
\item For $\mathbf{v}\in\mathcal{V}$ and $\psi\in\mathcal{W}'$, $\bigl(T''\Lambda_{\mathcal{V}}(\mathbf{v})\bigr)(\psi)=\Lambda_{\mathcal{V}}(\mathbf{v})(T'\psi)=(T'\psi)(\mathbf{v})=\psi(T\mathbf{v})=\Lambda_{\mathcal{W}}(T\mathbf{v})(\psi)$.\\
Hence $T''\Lambda_{\mathcal{V}}=\Lambda_{\mathcal{W}}T$ by Assumption \ref{assump:function-extensionality} on $\mathcal{W}'$ and $\mathcal{V}$.
\item If $\mathcal{V}$ and $\mathcal{W}$ are finite-dimensional, then $\Lambda_{\mathcal{V}}$ and $\Lambda_{\mathcal{W}}$ are isomorphisms by Theorem \ref{thm:algebraic-reflexivity-natural-map}(ii). Thus $T''=T''I_{\mathcal{V}''}=T''\Lambda_{\mathcal{V}}\Lambda_{\mathcal{V}}^{-1}=\Lambda_{\mathcal{W}}T\Lambda_{\mathcal{V}}^{-1}$ by (i) and Proposition \ref{prop:linear-map-product-algebra}(ii)--(iii).
\end{enumerate}
\hlt{Structural Commentary}: Pointwise evaluation proves naturality; finite reflexivity transports $T$ to its double dual.\\
\hlt{Key Insight}: Both routes evaluate each functional at the same image vector.
\end{proof}
\end{corollary}
